%% kapa: exam=1 round=32 year=2021 period=2 form=A subject=appraisal_law subject_ko=감정평가관계법규 seq=20 number=20
%% source: Q-Net 감정평가사 기출문제 (cst00309) artlSeq=5208397 file=extracted/1008/5208397/A형/2021년도 제32회 감정평가사 1차 2교시 A형.pdf
%% conversion: kordoc 4.24.51 -> tools/gen1.py (auto); review=proofread-sonnet; flags=none
\begin{question}{20}
국유재산법령상 국유재산에 관한 설명으로 옳지 않은 것은?
\begin{choices}
  \item 국유재산책임관의 임명은 중앙관서의 장이 소속 관서에 설치된 직위를 지정하는 것으로 갈음할 수 있다.
  \item 확정판결에 따라 일반재산에 대하여 사권(私權)을 설정할 수 있다.
  \item 총괄청은 국가에 기부하려는 재산이 재산가액 대비 유지·보수 비용이 지나치게 많은 경우에는 기부 받아서는 아니 된다.
  \item 국가 외의 자는 기부를 조건으로 하더라도 국유재산에 영구시설물을 축조할 수 없다.
  \item 중앙관서의 장은 국유재산의 관리·처분에 관련된 법령을 개정하려면 그 내용에 관하여 총괄청 및 감사원과 협의하여야 한다.
\end{choices}
\end{question}
\answer{4}
